% !TeX root = ../thuthesis-example.tex
\chapter{绪论}

\section{研究背景与意义}

\subsection{研究背景}
在数字化学习和办公环境中，时间管理、日程留痕和个人消费管理已经成为大学生和普通办公用户最常见的三类需求。当前市场上的相关应用虽然数量较多，但多数产品仅覆盖其中一种功能，例如番茄钟软件只提供倒计时服务，记账软件只关注收支录入，日记软件则更偏向文本记录，用户往往需要在多个应用之间频繁切换，导致信息分散、使用成本升高，也不利于形成连续、统一的个人行为数据。基于此，设计一套能够同时支持番茄计时、日记记录和基础记账分析的桌面应用系统具有较强的现实意义。

\subsection{研究意义}
本课题的研究意义主要体现在两个方面。其一，从应用价值角度看，该系统能够帮助用户在同一平台内完成专注计时、日记留存和消费管理，提高信息整合程度，增强使用连贯性。其二，从工程实践角度看，本课题以C++和Qt为核心技术路线，围绕界面设计、数据组织、模块协作、图表可视化和交互优化等内容展开实现，具有较强的软件工程训练价值。通过该项目，可以系统地验证桌面端轻量级个人信息管理软件的可行实现路径，为同类应用的开发提供参考。
从毕业设计的角度进一步分析，本系统兼具“功能完整性”和“技术可展示性”两方面特点。一方面，系统覆盖了用户登录、信息记录、数据统计、界面优化等多个完整业务流程，能够体现需求分析、系统设计、编码实现和测试验证的完整开发过程；另一方面，系统涉及Qt事件机制、定时器、文件系统、JSON序列化、自定义绘图、主题切换等多项关键技术，具备较强的综合训练意义。因此，本课题不仅适合用于实际场景中的个人效率管理，也能够较好地体现计算机科学与技术专业学生在桌面软件开发方向上的综合能力。

\section{国内外研究现状}

\subsection{国外研究现状}
国外在时间管理和个人效率工具方面起步较早，番茄工作法相关应用和任务管理工具已经较为成熟，一些产品还结合了统计图表、提醒机制和云同步能力，强调跨终端协同与用户行为数据分析。同时，国外桌面应用开发在跨平台框架和模块化设计方面积累了大量经验，Qt等技术方案被广泛应用于高性能桌面软件开发。

\subsection{国内研究现状}
国内相关软件产品在移动端发展迅速，围绕轻量记账、习惯打卡和时间管理形成了较为丰富的应用生态，但桌面端产品整体较少，且很多系统偏重商业化运营，存在功能堆叠、广告较多、数据导出受限等问题。对高校毕业设计而言，基于本地化存储与自主可控架构实现一套可独立运行的桌面系统，更具有研究与实践价值。

\subsection{现有问题分析}
综合现有产品可以发现，当前系统主要存在以下问题：一是功能分散，时间管理、日记与记账相互独立；二是部分系统对云端依赖程度较高，不利于隐私保护；三是桌面端产品在交互细节、图表展示和主题适配方面仍有提升空间。因此，本文选择以“集成化、轻量化、本地化”为目标，设计并实现番茄计时与记账日记系统。
此外，现有很多产品在“数据关联性”方面做得不够充分。例如，用户完成一次专注任务后，往往无法方便地将该事件与当天的日记内容和消费行为放在同一视角下观察，这使得个人行为数据之间缺乏关联。对于本项目而言，将专注记录、日记文本与账单信息纳入同一套本地管理体系，不仅可以提高软件实用性，也为后续进行更复杂的行为统计分析奠定了基础。

\sectionsection{研究内容与论文结构}

\subsection{研究内容}
本文围绕一个基于Qt的桌面系统展开研究与实现，主要内容包括：设计用户登录与会话机制；实现番茄计时和正向计时功能；构建日记记录和本地文件存储流程；新增记账模块并完成账单的录入、展示、编辑与删除；设计数据统计页面，通过饼状图和柱状图展示专注数据与消费分布；完成夜间主题、右键菜单和动画等交互优化；最后通过测试验证系统的可用性与稳定性。
在研究方法上，本文采用“需求驱动、逐步实现、测试验证”的开发路径。首先从用户日常使用场景出发，对系统所需功能进行抽象与拆分；其次根据Qt框架特点进行界面与逻辑分层实现；最后通过人工测试与QTest自动化测试相结合的方式，对关键功能进行验证。这样的研究过程既符合软件工程的开发规律，也便于将项目实现过程转化为较为完整的论文表达。

\subsection{论文结构安排}
本文共分为五章。第一章介绍研究背景、意义及国内外现状；第二章对系统需求和关键技术方案进行分析与比较；第三章重点阐述系统总体架构及各功能模块的设计与实现；第四章对系统进行测试并分析结果；第五章对全文进行总结，并给出后续改进方向。